% GENERATED FILE. Edit canonical lessons, then run npm run generate:books.
% canonical-source-hash: fnv1a64:92dbff361f25cea0
% canonical-lessons: HI-C538-sajona, HI-C538-taiyar, HI-C538-prabandh, HI-C538-sajna, HI-C538-binna

\chapter{Getting Things Ready}
\label{ch:hi-getting-things-ready}
\hlchaptermodality{538}

\begin{chapteropening}
\textbf{By the end of this chapter:} \emph{I can say how a feast or a celebration was got ready.}

Complete the last lesson of chapter 538: I can say how a feast or a celebration was got ready.
\end{chapteropening}

\section[\texorpdfstring{\hi{सँजोना} (sãjonā)}{सँजोना (sãjonā)}]{\hi{सँजोना} (sãjonā) — to treasure; to keep with care}

\label{lesson:HI-C538-sajona}



\begin{quote}
Before the new one: say the Hindi for an auspicious moment (to begin something), then the Hindi for to shift.
\end{quote}

\subsection*{You'll want to know: \hi{सँजोना}}
\textbf{\hi{सँजोना}} — \emph{sãjonā} — \textquotedblleft{}to treasure; to keep with care\textquotedblright{}.

\begin{quote}
\hi{दादी} \hi{ने} \hi{पुरानी} \hi{तस्वीरें} \hi{एक} \hi{डिब्बे} \hi{में} \hi{सँजोकर} \hi{रखी} \hi{हैं।}
\end{quote}

\begin{quote}
\emph{Dādī ne purānī tasvīreṁ ek ḍibbe meṁ sãjokar rakhī haiṁ.}
\end{quote}

\begin{quote}
Grandmother has kept the old photographs carefully in a box.
\end{quote}

Past, with \hi{ने} (he or she …): \hi{उसने} \hi{पुरानी} \hi{चिट्ठियाँ} \hi{सँजोईं}, \emph{usne purānī ciṭṭhiyā̃ sãjoīṁ}.

\begin{cousinweb}
\textbf{\hi{सँजोना}} (\emph{sãjonā}) is a tadbhava, inherited from Sanskrit \emph{saṁyōjayati}, \textquotedblleft{}brings together\textquotedblright{} (CDIAL 12990). CDIAL gives the Hindi verb as \textquotedblleft{}to decorate, make ready\textquotedblright{}; keeping something with care grew from that.
\end{cousinweb}

\begin{grammarlens}[title={What you've built}]
One more word for saying what happens if.
\end{grammarlens}

\subsection*{Your turn}
\begin{itemize}
\raggedright
  \item \emph{Say it:} \emph{sãjonā}
  \item \emph{Say it:} \emph{sãjonā}, once more
  \item \emph{Say it:} \emph{Dādī ne purānī tasvīreṁ ek ḍibbe meṁ sãjokar rakhī haiṁ.} — Grandmother has kept the old photographs carefully in a box.
  \item \emph{Recall:} say the Hindi for an auspicious moment (to begin something), then the Hindi for to shift, then say \emph{sãjonā} again
\end{itemize}

\begin{tcolorbox}[breakable,colback=teal!4,colframe=teal!35!black,title={Before you move on}]
What does \hi{सँजोना} mean? (\textquotedblleft{}To treasure; to keep with care\textquotedblright{}.) Say it once more.
\end{tcolorbox}
\section[\texorpdfstring{\hi{तैयार} (taiyār)}{तैयार (taiyār)}]{\hi{तैयार} (taiyār) — ready; prepared}

\label{lesson:HI-C538-taiyar}



\begin{quote}
Before the new one: say the Hindi for to shift, then the Hindi for to treasure.
\end{quote}

\subsection*{You'll want to know: \hi{तैयार}}
\textbf{\hi{तैयार}} — \emph{taiyār} — \textquotedblleft{}ready; prepared\textquotedblright{}.

\begin{quote}
\hi{खाना} \hi{तैयार} \hi{है}, \hi{सब} \hi{लोग} \hi{आइए।}
\end{quote}

\begin{quote}
\emph{Khānā taiyār hai, sab log āie.}
\end{quote}

\begin{quote}
The food is ready; come along, everyone.
\end{quote}

\begin{cousinweb}
\textbf{\hi{तैयार}} (\emph{taiyār}) is usually traced to Arabic \emph{ṭayyār}, \textquotedblleft{}swift, flying\textquotedblright{}, through Persian, where it came to mean \textquotedblleft{}ready\textquotedblright{}.
\end{cousinweb}

\begin{grammarlens}[title={What you've built}]
One more word for saying what happens if.
\end{grammarlens}

\subsection*{Your turn}
\begin{itemize}
\raggedright
  \item \emph{Say it:} \emph{taiyār}
  \item \emph{Say it:} \emph{taiyār}, once more
  \item \emph{Say it:} \emph{Khānā taiyār hai, sab log āie.} — The food is ready; come along, everyone.
  \item \emph{Recall:} say the Hindi for to shift, then the Hindi for to treasure, then say \emph{taiyār} again
\end{itemize}

\begin{tcolorbox}[breakable,colback=teal!4,colframe=teal!35!black,title={Before you move on}]
What does \hi{तैयार} mean? (\textquotedblleft{}Ready; prepared\textquotedblright{}.) Say it once more.
\end{tcolorbox}
\section[\texorpdfstring{\hi{प्रबंध} (prabandh)}{प्रबंध (prabandh)}]{\hi{प्रबंध} (prabandh) — management; arrangements}

\label{lesson:HI-C538-prabandh}



\begin{quote}
Before the new one: say the Hindi for to treasure, then the Hindi for ready.
\end{quote}

\subsection*{You'll want to know: \hi{प्रबंध}}
\textbf{\hi{प्रबंध}} — \emph{prabandh} — \textquotedblleft{}management; arrangements\textquotedblright{}.

\begin{quote}
\hi{दावत} \hi{का} \hi{प्रबंध} \hi{मेरी} \hi{चाची} \hi{ने} \hi{किया।}
\end{quote}

\begin{quote}
\emph{Dāvat kā prabandh merī cācī ne kiyā.}
\end{quote}

\begin{quote}
My aunt saw to the arrangements for the feast.
\end{quote}

Masculine.

\begin{cousinweb}
\textbf{\hi{प्रबंध}} (\emph{prabandh}) is a tatsama, taken whole from Sanskrit \emph{prabandha}, \textquotedblleft{}a connection, a tie\textquotedblright{} (Monier-Williams). Hindi uses it for managing and arranging.
\end{cousinweb}

\begin{grammarlens}[title={What you've built}]
One more word for saying what happens if.
\end{grammarlens}

\subsection*{Your turn}
\begin{itemize}
\raggedright
  \item \emph{Say it:} \emph{prabandh}
  \item \emph{Say it:} \emph{prabandh}, once more
  \item \emph{Say it:} \emph{Dāvat kā prabandh merī cācī ne kiyā.} — My aunt saw to the arrangements for the feast.
  \item \emph{Recall:} say the Hindi for to treasure, then the Hindi for ready, then say \emph{prabandh} again
\end{itemize}

\begin{tcolorbox}[breakable,colback=teal!4,colframe=teal!35!black,title={Before you move on}]
What does \hi{प्रबंध} mean? (\textquotedblleft{}Management; arrangements\textquotedblright{}.) Say it once more.
\end{tcolorbox}
\section[\texorpdfstring{\hi{सजना} (sajnā)}{सजना (sajnā)}]{\hi{सजना} (sajnā) — to be decorated; to be dressed up}

\label{lesson:HI-C538-sajna}



\begin{quote}
Before the new one: say the Hindi for ready, then the Hindi for management.
\end{quote}

\subsection*{You'll want to know: \hi{सजना}}
\textbf{\hi{सजना}} — \emph{sajnā} — \textquotedblleft{}to be decorated; to be dressed up\textquotedblright{}.

\begin{quote}
\hi{शादी} \hi{के} \hi{लिए} \hi{पूरा} \hi{घर} \hi{सज} \hi{रहा} \hi{है।}
\end{quote}

\begin{quote}
\emph{Śādī ke liye pūrā ghar saj rahā hai.}
\end{quote}

\begin{quote}
The whole house is being decorated for the wedding.
\end{quote}

Past (it …): \hi{वह} \hi{कमरा} \hi{सजा}, \emph{vah kamrā sajā}.

\begin{cousinweb}
\textbf{\hi{सजना}} (\emph{sajnā}) is probably related to Hindi \emph{sājnā}, \textquotedblleft{}to be prepared\textquotedblright{}, which CDIAL files under Sanskrit \emph{sajjyatē}, \textquotedblleft{}is equipped, is prepared\textquotedblright{} (CDIAL 13093); CDIAL does not list \emph{sajnā} itself; it is the intransitive partner of \emph{sajānā}, \textquotedblleft{}to decorate\textquotedblright{}.
\end{cousinweb}

\begin{grammarlens}[title={What you've built}]
One more word for saying what happens if.
\end{grammarlens}

\subsection*{Your turn}
\begin{itemize}
\raggedright
  \item \emph{Say it:} \emph{sajnā}
  \item \emph{Say it:} \emph{sajnā}, once more
  \item \emph{Say it:} \emph{Śādī ke liye pūrā ghar saj rahā hai.} — The whole house is being decorated for the wedding.
  \item \emph{Recall:} say the Hindi for ready, then the Hindi for management, then say \emph{sajnā} again
\end{itemize}

\begin{tcolorbox}[breakable,colback=teal!4,colframe=teal!35!black,title={Before you move on}]
What does \hi{सजना} mean? (\textquotedblleft{}To be decorated; to be dressed up\textquotedblright{}.) Say it once more.
\end{tcolorbox}
\section[\texorpdfstring{\hi{बीनना} (bīnnā)}{बीनना (bīnnā)}]{\hi{बीनना} (bīnnā) — to pick out, to pick over (grain)}

\label{lesson:HI-C538-binna}



\begin{quote}
Before the new one: say the Hindi for management, then the Hindi for to be decorated.
\end{quote}

\subsection*{You'll want to know: \hi{बीनना}}
\textbf{\hi{बीनना}} — \emph{bīnnā} — \textquotedblleft{}to pick out, to pick over (grain)\textquotedblright{}.

\begin{quote}
\hi{माँ} \hi{ने} \hi{दाल} \hi{बीनी।}
\end{quote}

\begin{quote}
\emph{Māṁ ne dāl bīnī.}
\end{quote}

\begin{quote}
Mother picked over the lentils.
\end{quote}

Past, with \hi{ने} (he or she …): \hi{उसने} \hi{चावल} \hi{बीने}, \emph{usne cāval bīne}.

\begin{cousinweb}
\textbf{\hi{बीनना}} (\emph{bīnnā}) is a tadbhava, inherited from Sanskrit \emph{vicinōti}, \textquotedblleft{}picks out, chooses\textquotedblright{} (CDIAL 11686). CDIAL gives the Hindi verb as \textquotedblleft{}to pick up, choose\textquotedblright{}.
\end{cousinweb}

\begin{grammarlens}[title={What you've built}]
One more word for saying what happens if.
\end{grammarlens}

\subsection*{Your turn}
\begin{itemize}
\raggedright
  \item \emph{Say it:} \emph{bīnnā}
  \item \emph{Say it:} \emph{bīnnā}, once more
  \item \emph{Say it:} \emph{Māṁ ne dāl bīnī.} — Mother picked over the lentils.
  \item \emph{Recall:} say the Hindi for management, then the Hindi for to be decorated, then say \emph{bīnnā} again
\end{itemize}

\begin{tcolorbox}[breakable,colback=teal!4,colframe=teal!35!black,title={Before you move on}]
What does \hi{बीनना} mean? (\textquotedblleft{}To pick out, to pick over (grain)\textquotedblright{}.) Say it once more.
\end{tcolorbox}
